\documentclass[a4paper,12pt]{jsarticle}
\usepackage[margin=25mm]{geometry}
\usepackage{amsmath,amssymb}
\usepackage{enumitem}

\begin{document}

\section*{第1問}
ある商品の需要を $x$（百個）とすると、1か月の利益 $P(x)$（万円）は
\[
P(x)=-2x^2+24x-40
\]
で表されるとする。ただし $0\le x\le 10$。

\begin{enumerate}[label=\arabic*.]
  \item 利益が最大となるときの $x$ を求めよ。
  \item 最大利益を求めよ。
  \item 利益が 0 以上となる $x$ の範囲を求めよ。
  \item 需要が $x=3$ から $x=7$ に増えるとき、利益は何万円増えるか。
\end{enumerate}

\section*{第2問}
文化祭の準備で、1から9までの番号が書かれたカードを使って抽選を行うことになった。  
カードは各番号1枚ずつあり、よく混ぜたあと同時に3枚を取り出す。取り出した3枚の番号は小さい順に $a<b<c$ として扱う。  
担当の生徒が「合計が偶数になることはどれくらい起こるか」「最大が7になる場面はどの程度あるか」などを調べることにした。

\begin{enumerate}[label=\arabic*.]
  \item 取り出し方は全部で何通りか。
  \item 3枚の和が偶数となる確率を求めよ。
  \item 3枚の最大値が $7$ である確率を求めよ。
  \item $a+b>c$ となる確率を求めよ。
\end{enumerate}

\section*{第3問}
三角形 $ABC$ について、辺の長さと角の情報を用いて形状を調べる。  
$AB=9,\ AC=7,\ \angle A=60^\circ$ とする。  
この三角形を土地の区画のモデルとみなし、辺 $BC$ を境界線、頂点 $A$ を観測点と考える。  
面積や高さなど、測量で必要になる量を順に求める。

\begin{enumerate}[label=\arabic*.]
  \item 辺 $BC$ の長さを求めよ。
  \item 三角形 $ABC$ の面積を求めよ。
  \item 外接円の半径 $R$ を求めよ。
  \item 辺 $BC$ を底辺としたときの高さを求めよ。
\end{enumerate}

\section*{第4問}
あるクラス40人の小テスト結果を分析する。  
集計の結果、平均点は $62$ 点、標準偏差は $10$ 点であった。  
得点を $x$ 点とするとき、偏差値 $H$ を
\[
H=50+10\frac{x-62}{10}
\]
で定める。  
先生は「基準点からどれだけ上か下か」を偏差値で確認し、学習状況を把握したいと考えている。

\begin{enumerate}[label=\arabic*.]
  \item $x=72$ のときの偏差値を求めよ。
  \item 偏差値が $60$ 以上となる得点 $x$ の範囲を求めよ。
  \item 得点が $52$ 点の生徒と $77$ 点の生徒の偏差値の差を求めよ。
  \item 偏差値 $45$ の生徒の得点を求めよ。
\end{enumerate}

\end{document}
